\input{../tex-inputs/latex-leseansicht-vorspann}

\section[Arthur Schnitzler an Gustav Schwarzkopf, 20. 7. 1904]{L04212 Arthur Schnitzler an Gustav Schwarzkopf, 20. 7. 1904}
\nopagebreak\mylabel{L04212v}
\rehead{ }\normalsize\beginnumbering\briefempfaengerindex{Schwarzkopf, Gustav@\textsc{Schwarzkopf, Gustav}!zzzSchnitzler, Arthur@\emph{von Arthur Schnitzler}!1904-07-201@{20. 7. 1904}|(be}
\toendnotes[C]{\smallbreak\pagebreak[2]}
\correspDesc{Versand  durch Arthur Schnitzler am 20. 7. 1904 in Reichenau an der Rax
\newline{}Zustellung  am 20. 7. 1904 in Wien
\newline{}Erhalt  durch Gustav Schwarzkopf im Zeitraum [20. 7. 1904 –
                  31. 8. 1904?] in Wien}\toendnotes[C]{\smallbreak}
\Standort{CUL, Schnitzler, B 96.}
\physDesc{Telegramm, 101 Zeichen
\newline{}Handschrift: lila Tinte, deutsche Kurrent
\newline{}Versand: 1) »\noindent{}\textcolor{gray}{\textbf{Eingelangt am}}{ }20/7 \textcolor{gray}{\textbf{190…}}{ }\textcolor{gray}{N}{ / }\textcolor{gray}{\textbf{Von}}{ }Reichenau NÖ.\oindex{Reichenau an der Rax@\textbf{Reichenau an der Rax}, \emph{Verwaltungsgebiet}|pw}{ / }\textcolor{gray}{\textbf{Aufgabe-Nr}}{ }46325 \textcolor{gray}{\textbf{mit … Taxworten
                                          (}}16\textcolor{gray}{\textbf{Worten … Chiffern)}}{ / }\textcolor{gray}{\textbf{Aufgegeben am … um}}{ }9\textsuperscript{\textcolor{gray}{1}00}«  2) Stempel: »\nobreak{}\oindex{I., Innere Stadt@\textbf{I., Innere Stadt}, \emph{Verwaltungsgebiet}|pwk}Wien 1/1, 20 VII 0{[}4{]}, 11–V\nobreak{}«. }\toendnotes[C]{\smallbreak}\pstart{}{\pb}\textsc{Gustav
                     Schwarzkopf}\pend{}\pstart{}\textsc{Tiefer
                        Graben 23} Wien\oindex{Wien@\textbf{Wien}!I., Innere Stadt@\textbf{I., Innere Stadt}!Tiefer Graben 23@\textbf{Tiefer Graben 23}, \emph{Wohngebäude}|pw}\pend{}{\bigskip}\vspace{1em}
\pstart
           \noindent{}{\pb}Wann \label{K_L04212-1v}\edtext{ko{\geminationm}en}{\lemma{\textnormal{\emph{kommen}}}\Cendnote{\textnormal{Schwarzkopf\pwindex{Schwarzkopf, Gustav 7.\,11.\,1853 Wien – 13.\,11.\,1939 ebd.@\textsc{Schwarzkopf, Gustav} (7.\,11.\,1853 Wien – 13.\,11.\,1939 ebd.), \emph{Schriftsteller}|pwk} hielt sich
                  in Sekirn\oindex{Sekirn@\textbf{Sekirn}|pwk} auf und kam nicht nach Reichenau\oindex{Reichenau an der Rax@\textbf{Reichenau an der Rax}, \emph{Verwaltungsgebiet}|pwk}.}}}\label{K_L04212-1} Sie endlich \strikeout{Wien} wir bleiben bis \label{K_L04212-2v}\edtext{Freitag}{\lemma{\textnormal{\emph{Freitag}}}\Cendnote{\textnormal{Die Abreise verzögerte sich auf Samstag, 23. 7. 1904.}}}\label{K_L04212-2}\pend
           \pstart Herzlichst \spacefill\mbox{Arthur–}\pend{}\selectlanguage{ngerman}\endnumbering\briefempfaengerindex{Schwarzkopf, Gustav@\textsc{Schwarzkopf, Gustav}!zzzSchnitzler, Arthur@\emph{von Arthur Schnitzler}!1904-07-201@{20. 7. 1904}|)be}\mylabel{L04212h}
\begin{anhang}
\end{anhang}\newcommand{\dateiname}{L04212}\newcommand{\titel}{Arthur Schnitzler an Gustav Schwarzkopf, 20. 7. 1904}\newcommand{\editorInnen}{Herausgegeben von Jahnke, SelmaMüller, Martin Anton}\input{../tex-inputs/latex-leseansicht-abspann}
